\uuid{DpHu}
\exo7id{4082}
\titre{exo7 4082}
\auteur{quercia}
\organisation{exo7}
\datecreate{2010-03-11}
\isIndication{false}
\isCorrection{true}
\chapitre{Equation différentielle}
\sousChapitre{Equations différentielles linéaires}
\cpge{mpsi-mp2i.primitives-equations-differentielles.a, pcsi.techniques-calcul-primitives.a, ptsi.techniques-calcul-primitives.a, mpsi-mp2i.primitives-equations-differentielles.b, pcsi.techniques-calcul-primitives.b, ptsi.techniques-calcul-primitives.b, mpsi-mp2i.integration.e, pcsi.integration.d, ptsi.integration.d, bcpst1.integration-segment.c, bcpst1.equations-differentielles-lineaires.a}
\module{Analyse}
\niveau{L2}
\difficulte{}

\contenu{
\question{
Trouver les fonctions $f,g$ continues vérifiant~:
$ \int_{t=0}^x f(t)\,d t = x - 1 + g(x)$ et
$ \int_{t=0}^x g(t)\,d t = x - 1 + f(x)$.
}
\reponse{
Poser $F(x) =  \int_{t=0}^x f(t)\,d t$ et $G(x) =  \int_{t=0}^x g(t)\,d t$
puis résoudre~: $$\left\{\begin{array}{lll} F(x) &= &x-1+G'(x)\\ G(x) &= &x-1 + F'(x)\\ F(0) &= &G(0) = 0.\hfill\\ \end{array}\right.$$

On trouve $f(x) = g(x) = 1$.
}
}
